\section{The Four Regimes}
\label{sec:four-regimes}

The complete experimental programme --- three polytope geometries, five
spectral-only baselines, two non-geometric controls, and an emanation
architecture --- reveals four distinct training outcomes. Together they form
a taxonomy of bridge behavior under the Steersman.

\subsection{Regime Classification}

\begin{table*}[t]
\centering
\small
\begin{tabular}{llcccp{4.5cm}}
\toprule
\textbf{Regime} & \textbf{Experiments} & \textbf{Fiedler} & \textbf{Co/Cross} & \textbf{Eigenvalues} & \textbf{Mechanism} \\
\midrule
\textbf{1: Block-Diagonal} & H-ch6, O-001, T-001r2 & $\sim$10\textsuperscript{-5}--10\textsuperscript{-4} & 10\textsuperscript{4}--10\textsuperscript{5}:1 & Clean $k{+}k$ split & Geometric contrastive: valid polytope pair spec \\
\textbf{2: Spectral Attractor} & H-ch3/4/8/12, E-001 & $0.084$--$0.102$ & $\sim$1:1 & Smooth distribution & Spectral-only or non-directional architecture \\
\textbf{3: Wrong-Partition} & WL-001 & $\sim$10\textsuperscript{-5} & $\sim$0 & $k{+}k$ split, wrong direction & Contrastive with incoherent pair spec \\
\textbf{4: Collapse} & R-001 & $\sim$10\textsuperscript{-5} & $\sim$0 & Chain-like progression & Non-geometric pair structure \\
\bottomrule
\end{tabular}
\caption{The four training regimes. Block-diagonal (Regime~1) requires all
three ingredients: cybernetic feedback, contrastive loss, and geometrically
coherent pair specification. The spectral attractor (Regime~2) is the
bridge's unprogrammed equilibrium. Wrong-partition (Regime~3) produces
eigenvalue structure without directional content. Collapse (Regime~4) produces
neither blocks nor connectivity.}
\label{tab:four-regimes-full}
\end{table*}

The regimes are distinguished by two independent diagnostics:

\begin{enumerate}
  \item \textbf{Eigenvalue structure:} Does the Laplacian spectrum exhibit a
    clean $k{+}k$ split (near-zero eigenvalues followed by large eigenvalues)?
    Regimes 1 and 3 exhibit this structure; Regimes 2 and 4 do not.
  \item \textbf{Directional coherence:} Does the co-planar/cross-planar
    coupling ratio indicate systematic directional preference? Only Regime~1
    produces this; all other regimes have co/cross $\sim$1 or $\sim$0.
\end{enumerate}

This creates a $2 \times 2$ classification:

\begin{table}[h]
\centering
\begin{tabular}{l|cc}
 & \textbf{Has eigenvalue structure} & \textbf{No structure} \\
\hline
\textbf{Directional} & Regime 1 (BD) & --- \\
\textbf{Non-directional} & Regime 3 (wrong-partition) & Regime 4 (collapse) \\
\end{tabular}
\caption{The $2 \times 2$ regime matrix. Regime~2 (spectral attractor)
precedes this classification --- it is the default state before any pair
specification is applied.}
\label{tab:regime-matrix}
\end{table}

Figure~\ref{fig:four-regimes} illustrates the four regimes with representative
trajectories.

\begin{figure*}[t]
  \centering
  \includegraphics[width=\textwidth]{paper4/figures/fig_four_regimes.pdf}
  \caption{The four regimes of topology programming. \textbf{(a)}~Block-diagonal:
  O-001 co/cross ratio climbs monotonically to 473{,}622:1.
  \textbf{(b)}~Spectral attractor: Fiedler converges to $\sim$0.09 regardless
  of channel count. \textbf{(c)}~Hierarchical coherence (emanation): Fiedler
  and co/cross both converge to the spectral attractor, confirming hierarchy
  provides connectivity without directionality. \textbf{(d)}~Connectivity
  collapse: WL-001 and R-001 produce identical trajectories approaching zero.}
  \label{fig:four-regimes}
\end{figure*}

The fourth cell (directional without eigenvalue structure) is unoccupied in
our experiments. This absence may reflect a geometric constraint: directional
coherence requires the channel space to partition into coupled blocks, which
necessarily produces eigenvalue structure. Whether this cell is geometrically
impossible or merely unvisited remains an open question.

\subsection{Regime Transitions}

The regimes form a hierarchy of requirements. Topology programming
(Regime~1) requires three simultaneous conditions:

\begin{enumerate}
  \item \textbf{Cybernetic feedback} (the Steersman): Without the
    sensor--oscilloscope--steersman--actuator loop, bridge matrices show no
    systematic structure development.
  \item \textbf{Contrastive loss}: Without pair-based attraction/repulsion,
    the bridge converges to the spectral attractor (Regime~2) regardless of
    channel count or architecture.
  \item \textbf{Geometrically coherent pair specification}: Without valid
    polytope geometry, contrastive loss produces either structure without
    direction (Regime~3) or complete collapse (Regime~4).
\end{enumerate}

Removing any single condition produces a predictable degradation:
\begin{itemize}
  \item Remove contrastive loss $\rightarrow$ Regime~2 (spectral attractor)
  \item Replace geometric pairs with random pairs $\rightarrow$ Regime~3
    (wrong-partition)
  \item Replace geometric pairs with algebraic pairs $\rightarrow$ Regime~4
    (collapse)
\end{itemize}

The distinction between Regimes~3 and~4 is diagnostic. Wrong-labels (random
partition) creates valid block structure because \emph{any} partition of $n$
channels into $k$ pairs induces a $k$-way bisection --- the contrastive loss
acts on the partition structure. But the resulting blocks have no internal
orientation: co-axial coupling equals cross-axial coupling within the blocks.
Resonance (number-theoretic) fails even to create clean blocks: the overlapping
pair structure (some channels appear in multiple relationships) produces a
connectivity \emph{gradient} rather than a partition.

\subsection{The Spectral Attractor as Ground State}

Regime~2 occupies a distinguished role: it is the bridge's \emph{ground
state}. Five independent experiments across $n \in \{3, 4, 6, 8, 12\}$ and
one architectural variant (emanation) converge to Fiedler $\approx 0.09$
(\S\ref{sec:spectral-attractor}). The contrastive loss does not add
connectivity to this ground state --- it \emph{redirects} it. The total
coupling budget remains approximately constant; the contrastive loss
concentrates it along the specified axes while suppressing it everywhere
else. This concentration produces the 1{,}130$\times$ Fiedler drop (from
0.09 to $9.0 \times 10^{-5}$) observed when comparing spectral-only to RD
contrastive at $n{=}6$.
